\subsection{Examination rewards a middling specificity}\label{app:registration}

Across 3.59 million debut filings (57.6\% base rate), completion is an
inverse-U in \emph{atypicality}: it climbs from 46.6\% at the least atypical
end to about 58\% in the middle and falls back to 53.4\% at the most atypical
(Figure~\ref{fig:profiles}, panel C). The same shape appears on the
past-facing level alone, more mildly, running 54.0\% at the bottom quintile to
58.4--59.1\% in the middle and 57.7\% at the top (Figure~\ref{fig:outcomes}A).

The signed axis behaves differently. Completion on $L$ falls to a trough of
51.9\% just above the median and then rises steadily to 58.0\% at the leading
end. That trough is composition rather than timing. Lead magnitude
correlates with atypicality at $+0.31$, so the filings near zero lead are
disproportionately the low-atypicality filings, and those complete poorly
for the reason above. Splitting by atypicality confirms it: the trough is
absent in the upper three fifths and survives only in the lowest. The
signed axis still enters a linear specification strongly
(harmonized estimates in the companion paper), but monotonically rather than through an interior
peak: what curvature exists at registration belongs to how far a filing's
language sits from its category, not to which side of it the filing sat.

\begin{figure}[h]\centering
\includegraphics[width=\textwidth]{results/debut_outcome_topic.png}
\caption{Outcomes for debut filings (owner's first-ever filing) across all
45 Nice classes under theme scoring, filing years 1995--2018, in quantile
bins with 95\% intervals, quantiles cut on the pooled distribution.
$n = 3{,}589{,}068$; registered subset $n = 2{,}066{,}600$. \emph{A:}
P(reached registration). \emph{B:} P(owner ever in SEC EDGAR, the
Commission's public filing system $\mid$ registration).}
\label{fig:outcomes}
\end{figure}

\begin{figure}[h]\centering
\includegraphics[width=\textwidth]{results/quintile_profiles.png}
\caption{Quintile profiles for the three outcomes. \emph{A:} the share of
registrations failing the five-year proof of continued use, by lead
quintile, adjusted within Nice class and registration year with the
drafting and filer controls of Table~\ref{tab:gate_lpm}; the dashed line
is the rate for all registrations. The profile rises gently and
monotonically through the fourth quintile and then steps up: about seven
tenths of the top-to-bottom gap is the single step from the fourth
quintile to the fifth, seven times the average step below it. \emph{B:}
the share of first-time filers whose owner ever reaches SEC reporting, by
quintile of atypicality and of lead, adjusted within Nice class and filing
year with description length held. \emph{C:} raw registration rates for
debut filings in twenty bins, cut on the pooled distribution --- an
inverse-U on atypicality, and on lead a trough just above the median that
is compositional, absent once atypicality is held in its middle fifth
(dashed). Horizontal line at the pooled rate.}
\label{fig:profiles}
\end{figure}

\begin{figure}[h]\centering
\includegraphics[width=0.9\textwidth]{results/fig_registration_sankey.png}
\caption{The registration step as a flow, unique serials filed 1995--2018
($n = 6{,}877{,}795$), split by counsel of record on the filing.
Represented applications reach registration at 61.9\%, self-filed at
45.6\%.}
\label{fig:sankey}
\end{figure}

The obvious objection is that unfamiliar language is simply refused more
often, so the measure recovers nothing beyond familiarity. Completion falls at
\emph{both} ends, so the most category-typical filings also complete less than
the middle, which familiarity cannot explain; and the curvature is specific to
the unsigned axis, the signed axis showing no interior peak at all.

Refiling after abandonment is itself mildly related to vocabulary position,
though not on the axis that matters here. Refiling rates are flat in
atypicality (a middle-minus-tails difference of $-0.06$pp on a 3.6\% base) but
rise with lead, from 2.6\% among the most lagging abandoned applications to
3.4\% among the most leading.

